\tmcapitulo{tmMorado}{Contar, ordenar, repetir y conectar: las matem\'aticas de lo que se puede enumerar.}{x}
\chapter{Matem\'atica discreta}
\label{cap:discreta}

\begin{tmsistemas}
Cu\'antas configuraciones tiene un sistema, cu\'antos casos de prueba hacen falta, cu\'an grande es el espacio de b\'usqueda de una clave y c\'omo se conectan los servidores de una red: todo eso se responde contando.
\end{tmsistemas}

\begin{center}
\begin{tabular}{|p{4.2cm}|p{4.8cm}|p{3.8cm}|}
\hline
\rowcolor{tmAzulClaro}\textbf{Tema} & \textbf{Qu\'e aprender\'as} & \textbf{D\'onde se usa}\\ \hline
2.1 Conteo & Sumar y multiplicar posibilidades & Pruebas, claves\\ \hline
2.2 Permutaciones y combinaciones & Orden s\'i u orden no & Selecci\'on de variables\\ \hline
2.3 Recurrencias & Costo de algoritmos & Divide y vencer\'as\\ \hline
2.4 Primos y divisibilidad & Factorizar y MCD & Criptograf\'ia\\ \hline
2.5 Aritm\'etica modular & Aritm\'etica con residuos & Contadores, RSA\\ \hline
2.6 Grafos & Redes y dependencias & Redes, IA, compiladores\\ \hline
\end{tabular}
\end{center}

% =====================================================================
\section{Principios de conteo}
\label{sec:2-1}
\tmgancho{Antes de probar todas las configuraciones de un sistema, \textquestiondown sabes cu\'antas son?}

\begin{tmconcepto}{sumar o multiplicar}
Si haces cosas una \textbf{despu\'es} de otra, \textbf{multiplicas}. Si eliges una cosa \textbf{o} la otra, \textbf{sumas}.
\end{tmconcepto}

\tmcomo{contar configuraciones}{{Identifica las etapas del procedimiento},{Cuenta las opciones de cada etapa},{\textquestiondown Van en secuencia? Multiplica},{Descarta las combinaciones no v\'alidas}}

\subsection{Principio aditivo}
\begin{tmdef}{Principio aditivo}
Si una tarea se hace de una forma \textbf{o} de otra, mutuamente excluyentes, el total es la suma: $N=n_1+n_2+\cdots+n_k$.
\end{tmdef}

\begin{tmejemplo}{Interfaz de comunicaci\'on}
Hay 2 opciones de conexi\'on serial y 3 inal\'ambricas. Se elige exactamente una: $N=2+3=5$ opciones.
\end{tmejemplo}

\subsection{Principio multiplicativo}
\begin{tmdef}{Principio multiplicativo}
Si un procedimiento tiene varias etapas y cada una admite $n_i$ posibilidades, el total es el producto: $N=n_1\cdot n_2\cdots n_k$.
\end{tmdef}

\begin{center}
\begin{tikzpicture}[level distance=1.5cm,every node/.style={font=\sffamily\small},
  level 1/.style={sibling distance=4.6cm},level 2/.style={sibling distance=2.2cm},
  tn/.style={draw=tmAzul,fill=tmAzulClaro,rounded corners=3pt,inner sep=3pt},
  hoja/.style={draw=tmVerde,fill=tmVerdeClaro,rounded corners=3pt,inner sep=3pt}]
  \node[tn] {inicio}
    child {node[tn] {Chrome}
      child {node[hoja] {Windows}}
      child {node[hoja] {Linux}}}
    child {node[tn] {Firefox}
      child {node[hoja] {Windows}}
      child {node[hoja] {Linux}}};
  \node[font=\sffamily\bfseries,text=tmVerde] at (6,-3) {$2\times2=4$ caminos};
\end{tikzpicture}
\end{center}

\begin{tmejemplo}{Configuraciones de prueba}
Un formulario web: 3 navegadores, 2 sistemas operativos y 4 perfiles de usuario, todas las combinaciones compatibles.
\begin{tmpasos}
\item Tres etapas independientes: $3$, $2$ y $4$ opciones.
\item $N=3\cdot2\cdot4=24$ configuraciones.
\item Es el tama\~no del espacio de pruebas, no el n\'umero de pruebas que har\'as: algunas ser\'an redundantes.
\end{tmpasos}
\end{tmejemplo}

\begin{tmejemplo}{Espacio de b\'usqueda de una clave}
Clave de 4 d\'igitos decimales, con repeticiones y ceros iniciales: $N=10^4=10\,000$ claves posibles.
\end{tmejemplo}

\begin{tmdef}{Principio del palomar}
Si se reparten $n+1$ objetos en $n$ cajas, alguna caja recibe al menos 2 objetos.
\end{tmdef}

\begin{tmsistemas}
El conteo dimensiona el espacio de b\'usqueda de una clave, pero \textbf{no} mide la seguridad real: el tiempo de un ataque depende de la velocidad de verificaci\'on, los bloqueos y c\'omo se proteja el sistema. El palomar explica por qu\'e las funciones hash \textbf{siempre} tienen colisiones.
\end{tmsistemas}

\begin{tmcuidado}Si ciertos navegadores no son compatibles con algunos sistemas operativos, las etapas dejan de ser independientes y el producto directo \textbf{sobrecuenta}. Hay que contar solo las combinaciones v\'alidas.\end{tmcuidado}

\begin{tmejercicios}[Ejercicios 2.1 \hfill {\small\mdseries respuestas en p\'ag. \pageref{resp:2-1}}]
\label{ej:2-1}
\begin{tmlista}
\item \tmbasico\ Un sistema usa 4 m\'etodos de autenticaci\'on y 3 niveles de acceso. \textquestiondown Cu\'antas configuraciones hay?
\item \tmbasico\ Una app se prueba en 4 navegadores, 3 sistemas operativos y 2 tama\~nos de pantalla. \textquestiondown Cu\'antas combinaciones?
\item \tmbasico\ \textquestiondown Cu\'antos PIN de 6 d\'igitos existen?
\item \tmmedio\ Una placa lleva 3 letras (26 posibles) y 3 d\'igitos. \textquestiondown Cu\'antas placas hay?
\item \tmmedio\ Debes elegir un servidor: 5 locales o 3 en la nube. \textquestiondown Cu\'antas opciones?
\item \tmmedio\ Una contrase\~na de 8 caracteres con 62 s\'imbolos posibles (con repetici\'on). \textquestiondown Cu\'antas hay?
\item \tmdificil\ En el ejercicio 2, Safari solo funciona en un sistema operativo (de los 3). \textquestiondown Cu\'antas combinaciones v\'alidas quedan?
\item \tmdificil\ Guardas 25 archivos en 24 carpetas. Justifica que alguna carpeta tiene al menos 2.
\end{tmlista}
\end{tmejercicios}

\tmreto{2}{Misi\'on 1 (Espacio de Pruebas): configura el espacio de pruebas hasta tener exactamente 24 combinaciones.}

\begin{tmresumen}
\begin{itemize}
\item Aditivo: ``o'' suma. Multiplicativo: ``y'' multiplica.
\item El conteo dimensiona espacios de b\'usqueda; no mide por s\'i solo la seguridad.
\item Si las etapas no son independientes, cuenta solo las combinaciones v\'alidas.
\end{itemize}
\end{tmresumen}

% =====================================================================
\section{Permutaciones y combinaciones}
\label{sec:2-2}
\tmgancho{\textquestiondown De cu\'antas formas se pueden ordenar tres tareas de un servidor? \textquestiondown Y cu\'antos grupos de revisi\'on puedes armar?}

\begin{tmconcepto}{orden s\'i u orden no}
\textbf{Permutaci\'on:} el orden importa (como un podio). \textbf{Combinaci\'on:} el orden no importa (como un equipo).
\end{tmconcepto}

\tmcomo{elegir la f\'ormula correcta}{{\textquestiondown Se repiten elementos? Usa $n^r$},{\textquestiondown Importa el orden? S\'i: $P(n, r)$},{Si no importa: $\binom nr$},{Verifica con un caso peque\~no}}

\subsection{Factorial}
\begin{tmdef}{Factorial}
$n!=n(n-1)(n-2)\cdots2\cdot1$ y por definici\'on $0!=1$. Cuenta las formas de ordenar $n$ objetos distintos: $4!=24$.
\end{tmdef}

\subsection{Cu\'ando importa el orden}
\begin{center}
\begin{tikzpicture}[every node/.style={font=\sffamily\small,rounded corners=4pt,inner sep=6pt,align=center}]
  \node[draw=tmAzul,fill=tmAzulClaro,text width=4.6cm] (a) at (0,0) {\textbf{Permutaci\'on}\\el orden \textbf{s\'i} importa\\$P(n, r)=\dfrac{n!}{(n-r)!}$};
  \node[draw=tmVerde,fill=tmVerdeClaro,text width=4.6cm] (b) at (6.2,0) {\textbf{Combinaci\'on}\\el orden \textbf{no} importa\\$\dbinom nr=\dfrac{n!}{r!(n-r)!}$};
  \node[font=\sffamily\scriptsize,text=tmGris] at (3.1,-1.7) {relaci\'on: $P(n, r)=r!\,\binom nr$};
\end{tikzpicture}
\end{center}

\begin{tmejemplo}{Planificaci\'on de tareas}
Se ejecutan 3 tareas distintas elegidas de 5. \textquestiondown Cu\'antos \'ordenes hay?
\begin{tmpasos}
\item El orden importa: $P(5, 3)=\dfrac{5!}{2!}=\dfrac{120}{2}=60$.
\item En la pr\'actica no todos son v\'alidos si hay dependencias entre tareas.
\end{tmpasos}
\end{tmejemplo}

\begin{tmejemplo}{Revisi\'on de m\'odulos}
De 5 m\'odulos se eligen 3 para una revisi\'on conjunta; solo importa cu\'ales:
\[\binom53=\frac{5!}{3!\,2!}=10.\]
\end{tmejemplo}

\begin{center}
\begin{tikzpicture}[every node/.style={font=\sffamily\small,circle,fill=tmAzulClaro,draw=tmAzul,inner sep=2pt,minimum size=0.7cm}]
  \foreach \n in {0,1,2,3,4,5}{
    \foreach \k in {0,...,\n}{
      \pgfmathtruncatemacro{\v}{int(round(exp(ln(factorial(\n))-ln(factorial(\k))-ln(factorial(\n-\k)))))}
      \node at ({(\k-\n/2)*0.95},{-\n*0.85}) {\v};}}
  \node[rectangle,draw=none,fill=none,anchor=west,text=tmGris,align=left] at (3.1,-2) {tri\'angulo de Pascal:\\$\binom nr=\binom{n-1}{r-1}+\binom{n-1}{r}$};
\end{tikzpicture}
\end{center}

\begin{tmtip}Pregunta clave: \textquestiondown cambia el resultado si intercambio dos elementos? Si s\'i, es permutaci\'on; si no, combinaci\'on. Siempre $\binom nr=\binom n{n-r}$.\end{tmtip}

\begin{tmsistemas}
\textbf{IA:} con 8 caracter\'isticas y 3 a elegir hay $\binom83=56$ subconjuntos para evaluar. \textbf{Ciberseguridad:} una contrase\~na de 6 letras min\'usculas tiene $26^6=308\,915\,776$ posibilidades. Ese n\'umero no dice qu\'e tan segura es: las personas no eligen claves al azar.
\end{tmsistemas}

\begin{tmcuidado}Una clave \textbf{con repetici\'on} se cuenta con $n^r$, no con $P(n, r)$. Permutaci\'on y combinaci\'on suponen elementos \textbf{sin repetir}.\end{tmcuidado}

\begin{tmejercicios}[Ejercicios 2.2 \hfill {\small\mdseries respuestas en p\'ag. \pageref{resp:2-2}}]
\label{ej:2-2}
\begin{tmlista}
\item \tmbasico\ \textquestiondown De cu\'antas formas se ordenan 4 tareas diferentes?
\item \tmbasico\ \textquestiondown Cu\'antos grupos de 4 se forman con 10 estudiantes?
\item \tmbasico\ De un sistema de 12 m\'odulos se eligen 5. \textquestiondown Cu\'antos subconjuntos hay?
\item \tmmedio\ Calcula $P(6, 2)$.
\item \tmmedio\ Con 8 estudiantes: \textquestiondown cu\'antas ternas con presidente, secretario y tesorero? \textquestiondown Y cu\'antos comit\'es de 3 sin cargos?
\item \tmmedio\ De 10 m\'odulos, elegir 4 para revisi\'on y, aparte, ordenar 4 para ejecuci\'on. Calcula ambos.
\item \tmdificil\ \textquestiondown Cu\'antas palabras distintas salen reordenando las letras de REDES?
\item \tmdificil\ Demuestra $\binom nr=\binom n{n-r}$ y calcula $\binom{10}3-\binom{10}7$.
\end{tmlista}
\end{tmejercicios}

\tmreto{2}{Contar es el primer paso del laboratorio del cap\'itulo 2: practica con el espacio de pruebas antes de pasar a las siguientes misiones.}

\begin{tmresumen}
\begin{itemize}
\item $P(n, r)=\frac{n!}{(n-r)!}$ si importa el orden; $\binom nr=\frac{n!}{r!(n-r)!}$ si no importa.
\item Con repetici\'on se usa $n^r$.
\item El resultado matem\'atico no sustituye el an\'alisis de seguridad real.
\end{itemize}
\end{tmresumen}

% =====================================================================
\section{Sucesiones y recurrencias}
\label{sec:2-3}
\tmgancho{\textquestiondown C\'omo sabe un ingeniero cu\'anto tardar\'a un algoritmo sin ejecutarlo para todos los tama\~nos?}

\begin{tmconcepto}{la recurrencia como ficha de domin\'o}
Una recurrencia es una receta que dice c\'omo obtener el siguiente t\'ermino a partir de los anteriores, como una cadena de domin\'os.
\end{tmconcepto}

\tmcomo{resolver una recurrencia simple}{{Escribe los primeros t\'erminos},{Busca el patr\'on},{Propone una f\'ormula cerrada},{Comprueba con $n=1, 2, 3$}}

\subsection{Sucesiones}
\begin{tmdef}{Sucesi\'on}
Lista ordenada de valores que sigue una regla. Aritm\'etica: $a_n=a_1+(n-1)d$. Geom\'etrica: $a_n=a_1r^{n-1}$. Ejemplo: $2, 4, 6, 8, 10,\dots$ con $a_n=2n$.
\end{tmdef}

\begin{tmdef}{Relaci\'on de recurrencia}
Define un t\'ermino a partir de los anteriores. Ejemplo: $a_n=a_{n-1}+3$ con $a_1=2$ da $2, 5, 8, 11, 14,\dots$
\end{tmdef}

\begin{tmejemplo}{Suma aritm\'etica}
Suma de los primeros 10 t\'erminos de $a_n=2n+3$.
\begin{tmpasos}
\item $a_1=5$ y $a_{10}=23$.
\item $S=\dfrac{n(a_1+a_n)}{2}=\dfrac{10\cdot28}{2}=140$.
\end{tmpasos}
\end{tmejemplo}

\begin{tmejemplo}{Fibonacci y torres de Hanoi}
\begin{tmpasos}
\item Fibonacci: $F_n=F_{n-1}+F_{n-2}$ da $1, 1, 2, 3, 5, 8, 13, 21,\dots$
\item Hanoi: mover $n$ discos exige $T(n)=2T(n-1)+1$ movimientos, con $T(1)=1$.
\item Resulta $T(n)=2^n-1$: con 5 discos son 31 movimientos.
\end{tmpasos}
\end{tmejemplo}

\subsection{An\'alisis de algoritmos: divide y vencer\'as}
Un algoritmo parte el problema en 2 mitades y combina en tiempo proporcional a $n$:
\tmformula{T(n)=2\,T(n/2)+cn\ \Rightarrow\ T(n)=\Theta(n\log n)}

\begin{center}
\begin{tikzpicture}[every node/.style={font=\sffamily\scriptsize,draw=tmMorado,fill=tmMorado!12,rounded corners=2pt,inner sep=2pt},
  level distance=1.1cm,level 1/.style={sibling distance=4cm},level 2/.style={sibling distance=2cm},level 3/.style={sibling distance=1cm}]
  \node {$cn$}
    child {node {$cn/2$} child {node {$cn/4$} child {node {}} child {node {}}} child {node {$cn/4$} child {node {}} child {node {}}}}
    child {node {$cn/2$} child {node {$cn/4$} child {node {}} child {node {}}} child {node {$cn/4$} child {node {}} child {node {}}}};
  \node[draw=none,fill=none,anchor=west,text=tmGris,font=\sffamily\small,align=left] at (4.7,-1.5) {cada nivel trabaja $cn$\\hay $\log_2 n$ niveles\\total: $cn\log_2 n$};
\end{tikzpicture}
\end{center}

\begin{tmsistemas}
Este patr\'on es el de \textbf{MergeSort}. En una estructura jer\'arquica, como un \'arbol binario completo, el nivel $k$ (con la ra\'iz en el nivel 0) tiene como m\'aximo $N_k=2^k$ nodos: el nivel 5 puede tener 32. En hardware, los contadores y las secuencias de estados son sucesiones discretas.
\end{tmsistemas}

\begin{tmtip}Para comparar dos algoritmos no los ejecutes: resuelve su recurrencia. $T(n)=2T(n/2)+n$ da $n\log n$; $T(n)=T(n-1)+n$ da $n^2/2$.\end{tmtip}

\begin{tmejercicios}[Ejercicios 2.3 \hfill {\small\mdseries respuestas en p\'ag. \pageref{resp:2-3}}]
\label{ej:2-3}
\begin{tmlista}
\item \tmbasico\ Escribe los primeros 6 t\'erminos de $a_n=3a_{n-1}$ con $a_1=2$.
\item \tmbasico\ Un algoritmo hace 4 operaciones por elemento de una lista de $n$. Expr\'esalo en funci\'on de $n$.
\item \tmbasico\ \textquestiondown Cu\'antos nodos puede tener el nivel 5 de un \'arbol binario completo?
\item \tmmedio\ Suma los primeros 10 t\'erminos de $a_n=2n+3$.
\item \tmmedio\ Escribe los 8 primeros n\'umeros de Fibonacci.
\item \tmmedio\ \textquestiondown Cu\'antos movimientos requiere Hanoi con 5 discos?
\item \tmdificil\ Resuelve $T(n)=T(n-1)+n$ con $T(1)=1$ y calcula $T(4)$.
\item \tmdificil\ Explica por qu\'e resolver una recurrencia ayuda a comparar dos algoritmos.
\end{tmlista}
\end{tmejercicios}

\tmreto{2}{Misi\'on 2 (Contador Modular): un contador cuya actualizaci\'on es una recurrencia modular. Usa lo aprendido aqu\'i para predecir su estado.}

\begin{tmresumen}
\begin{itemize}
\item Una sucesi\'on sigue una regla; una recurrencia define cada t\'ermino con los anteriores.
\item $T(n)=2T(n/2)+cn$ es $\Theta(n\log n)$; Hanoi es $2^n-1$.
\item Resolver recurrencias compara algoritmos sin ejecutarlos.
\end{itemize}
\end{tmresumen}

% =====================================================================
\section{N\'umeros primos y divisibilidad}
\label{sec:2-4}
\tmgancho{\textquestiondown Por qu\'e multiplicar dos n\'umeros primos es f\'acil, pero separarlos de nuevo es tan dif\'icil que protege tus mensajes?}

\begin{tmconcepto}{los primos como ladrillos}
Los primos son los \textbf{ladrillos} de los n\'umeros: todo entero mayor que 1 se arma multiplicando primos, y de una \'unica manera.
\end{tmconcepto}

\tmcomo{aplicar el algoritmo de Euclides}{{Divide $a$ entre $b$},{Toma el residuo $r$},{Repite con el par $(b, r)$},{Cuando $r=0$, el MCD es el \'ultimo divisor}}

\subsection{Conceptos}
\begin{tmdef}{Divisibilidad, primo y compuesto}
$a$ es divisible por $b>0$ si existe un entero $k$ con $a=bk$ (24 es divisible por 6: $24=6\cdot4$). Un \textbf{primo} es un entero mayor que 1 con exactamente dos divisores positivos: 2, 3, 5, 7, 11, 13. Un \textbf{compuesto} tiene m\'as de dos.
\end{tmdef}

\begin{tmdef}{Teorema fundamental}
Todo entero mayor que 1 se escribe de forma \'unica como producto de primos. Ejemplo: $60=2^2\cdot3\cdot5$.
\end{tmdef}

\begin{center}
\begin{tikzpicture}[every node/.style={font=\sffamily\small,circle,draw=tmAzul,fill=tmAzulClaro,inner sep=3pt,minimum size=0.8cm},level distance=1.2cm,
  level 1/.style={sibling distance=2.6cm},level 2/.style={sibling distance=1.3cm}]
  \node {60}
    child {node {2}}
    child {node {30}
      child {node[fill=tmVerdeClaro,draw=tmVerde] {2}}
      child {node {15} child {node[fill=tmVerdeClaro,draw=tmVerde] {3}} child {node[fill=tmVerdeClaro,draw=tmVerde] {5}}}};
\end{tikzpicture}
\end{center}

\begin{tmejemplo}{\textquestiondown 91 es primo?}
\begin{tmpasos}
\item Basta probar primos hasta $\sqrt{91}\approx9{,}5$: 2, 3, 5, 7.
\item $91=7\cdot13$. Es \textbf{compuesto}.
\end{tmpasos}
\end{tmejemplo}

\subsection{M\'aximo com\'un divisor y algoritmo de Euclides}
El $\operatorname{MCD}$ es el mayor entero que divide a ambos. Euclides: se divide y se repite con el residuo hasta llegar a 0.

\begin{center}
\begin{tabular}{|c|c|}
\hline
\rowcolor{tmAzulClaro}\textbf{Divisi\'on} & \textbf{Paso}\\ \hline
$48=2(18)+12$ & primera divisi\'on\\ \hline
$18=1(12)+6$ & segunda divisi\'on\\ \hline
$12=2(6)+0$ & residuo cero: $\operatorname{MCD}(48, 18)=6$\\ \hline
\end{tabular}
\end{center}

\begin{tmtip}Relaci\'on \'util: $\operatorname{MCD}(a, b)\cdot\operatorname{MCM}(a, b)=a\cdot b$. Con 12 y 18: $6\cdot36=216$.\end{tmtip}

\begin{tmcodigo}
def mcd(a, b):
    while b != 0:
        a, b = b, a % b
    return a

def es_primo(n):
    if n < 2:
        return False
    d = 2
    while d * d <= n:
        if n % d == 0:
            return False
        d += 1
    return True

print(mcd(48, 18), es_primo(91))   # 6 False
\end{tmcodigo}

\begin{tmsistemas}
En \textbf{RSA} se eligen dos primos distintos $p$ y $q$ y se calcula $n=pq$. Con $p=5$ y $q=11$ resulta $n=55$: sirve para ilustrar, pero \textbf{es demasiado peque\~no} para proteger informaci\'on real. Los sistemas reales exigen primos enormes, buena implementaci\'on y claves del tama\~no que fijan los est\'andares vigentes.
\end{tmsistemas}

\begin{tmcuidado}El 1 \textbf{no es primo} (solo tiene un divisor positivo) y el 2 es el \'unico primo par.\end{tmcuidado}

\begin{tmejercicios}[Ejercicios 2.4 \hfill {\small\mdseries respuestas en p\'ag. \pageref{resp:2-4}}]
\label{ej:2-4}
\begin{tmlista}
\item \tmbasico\ \textquestiondown 91 es primo o compuesto?
\item \tmbasico\ Factoriza 180 en primos.
\item \tmbasico\ Calcula $\operatorname{MCD}(84, 30)$ con Euclides.
\item \tmmedio\ Calcula $\operatorname{MCM}(12, 18)$.
\item \tmmedio\ \textquestiondown 97 es primo?
\item \tmmedio\ Calcula $\operatorname{MCD}(252, 105)$ con Euclides.
\item \tmdificil\ Explica por qu\'e basta probar divisores hasta $\sqrt n$.
\item \tmdificil\ Explica por qu\'e un ejemplo de RSA con primos peque\~nos no protege informaci\'on real.
\end{tmlista}
\end{tmejercicios}

\tmreto{2}{Pon a prueba tu lectura de n\'umeros primos y factorizaci\'on en las misiones del laboratorio del cap\'itulo 2.}

\begin{tmresumen}
\begin{itemize}
\item Todo entero $>1$ es producto \'unico de primos.
\item Euclides: $\operatorname{MCD}(a, b)=\operatorname{MCD}(b, a\bmod b)$.
\item RSA se apoya en la dificultad de factorizar $n=pq$ cuando $p$ y $q$ son enormes.
\end{itemize}
\end{tmresumen}

% =====================================================================
\section{Aritm\'etica modular}
\label{sec:2-5}
\tmgancho{\textquestiondown Qu\'e hora ser\'a dentro de 30 horas? Tu reloj ya sabe aritm\'etica modular.}

\begin{tmconcepto}{n\'umeros que dan vueltas}
En aritm\'etica modular los n\'umeros \textbf{dan vueltas} como las horas de un reloj: al llegar al m\'odulo, vuelves a empezar en 0.
\end{tmconcepto}

\tmcomo{calcular $a\bmod m$}{{Divide $a$ entre $m$},{Qu\'edate con el residuo},{En sumas y productos, reduce en cada paso},{El resultado est\'a entre $0$ y $m-1$}}

\tmformulanota{a\equiv b\pmod m}{$a$ y $b$ dejan el mismo residuo al dividir}{$m$ es el m\'odulo: el tama\~no del reloj}

\subsection{Congruencias}
\begin{tmdef}{Congruencia}
$a\equiv b\pmod m$ si $a$ y $b$ dejan el mismo residuo al dividirse entre $m$, es decir, si $m\mid(a-b)$. Ejemplo: $17\equiv5\pmod{12}$.
\end{tmdef}

\begin{center}
\begin{tikzpicture}[scale=0.95]
  \draw[tmAzul,line width=1.2pt] (0,0) circle (1.9);
  \foreach \i in {0,...,11}{\node[font=\sffamily\small\bfseries] at ({90-30*\i}:1.55) {\i};}
  \draw[-{Stealth},tmCoral,line width=1.4pt] (0,0) -- ({90-30*5}:1.05);
  \node[font=\sffamily\small,text=tmGris,anchor=west,align=left] at (2.4,0) {$17\equiv5\pmod{12}$\\$(8+9)\bmod7=3$ \quad $(4\cdot6)\bmod7=3$};
\end{tikzpicture}
\end{center}

\begin{tmejemplo}{Operaciones modulares}
\begin{tmpasos}
\item $(8+9)\bmod7=17\bmod7=3$.
\item $(4\cdot6)\bmod7=24\bmod7=3$.
\item $29\bmod6=5$ y $(15+23)\bmod8=38\bmod8=6$.
\end{tmpasos}
\end{tmejemplo}

\begin{tmejemplo}{Contador de m\'odulo 8}
Estados $0, 1,\dots,7$ y luego otra vez $0$: $x_{\text{nuevo}}=(x+1)\bmod8$. Si $x=7$, el siguiente es $(7+1)\bmod8=0$.
\end{tmejemplo}

\subsection{Inverso modular y RSA en miniatura}
El inverso de $a$ m\'odulo $m$ es $b$ con $ab\equiv1\pmod m$. Existe si $\operatorname{MCD}(a, m)=1$. Ejemplo: $3\cdot5=15\equiv1\pmod7$, as\'i que $3^{-1}=5$.

\begin{tmejemplo}{RSA de juguete}
$p=5$, $q=11$, $n=55$, $\varphi(n)=4\cdot10=40$. Se elige $e=3$ y $d=27$, pues $3\cdot27=81\equiv1\pmod{40}$.
\begin{tmpasos}
\item Mensaje $m=7$. Cifrado: $c=7^3\bmod55=343\bmod55=13$.
\item Descifrado: $13^{27}\bmod55=7$. Se recupera el mensaje.
\item Clave p\'ublica $(55, 3)$; clave privada $27$.
\end{tmpasos}
\end{tmejemplo}

\begin{tmsistemas}
\textbf{\'Indices circulares:} una cola circular de 10 posiciones en la posici\'on 8 que avanza 5 termina en $(8+5)\bmod10=3$. \textbf{Hash tables:} la posici\'on de un dato es $h(k)\bmod m$. \textbf{Blockchain:} usa criptograf\'ia (hash, firmas), pero tambi\'en estructuras de datos, consenso y redes; la aritm\'etica modular es solo una pieza.
\end{tmsistemas}

\begin{tmcodigo}
print(29 % 6)                  # 5
print((15 + 23) % 8)           # 6
print(pow(7, 3, 55))           # 13  (cifrado RSA de juguete)
print(pow(13, 27, 55))         # 7   (descifrado)
print(pow(3, -1, 7))           # 5   (inverso modular)
\end{tmcodigo}

\begin{tmcuidado}En Python \texttt{\%} con negativos da resultado no negativo: \texttt{-1 \% 5} es 4, justo lo que pide la matem\'atica. En otros lenguajes puede dar $-1$.\end{tmcuidado}

\begin{tmejercicios}[Ejercicios 2.5 \hfill {\small\mdseries respuestas en p\'ag. \pageref{resp:2-5}}]
\label{ej:2-5}
\begin{tmlista}
\item \tmbasico\ Calcula $29\bmod6$.
\item \tmbasico\ Calcula $(15+23)\bmod8$.
\item \tmbasico\ Un \'indice circular de 10 posiciones est\'a en la 8 y avanza 5. \textquestiondown D\'onde termina?
\item \tmmedio\ Calcula $(7\cdot9)\bmod5$.
\item \tmmedio\ Un contador de m\'odulo 16 est\'a en 13 y recibe 7 incrementos. \textquestiondown Qu\'e estado alcanza?
\item \tmmedio\ Halla el inverso de 3 m\'odulo 7.
\item \tmdificil\ Cifra $m=7$ con la clave p\'ublica $(55, 3)$.
\item \tmdificil\ Explica qu\'e aporta la aritm\'etica modular a RSA y por qu\'e no basta para describir toda una blockchain.
\end{tmlista}
\end{tmejercicios}

\tmreto{2}{Misi\'on 2 (Contador Modular): calcula el estado de un contador modular que parte de 13 con saltos de 7 en m\'odulo 16.}

\begin{tmresumen}
\begin{itemize}
\item $a\equiv b\pmod m$ si dejan el mismo residuo.
\item Se puede reducir cada paso m\'odulo $m$ en sumas y productos.
\item RSA cifra con $c=m^e\bmod n$ y descifra con $m=c^d\bmod n$.
\end{itemize}
\end{tmresumen}

% =====================================================================
\section{Introducci\'on a grafos}
\label{sec:2-6}
\tmgancho{\textquestiondown C\'omo sabe tu ruta de Internet cu\'al es el camino m\'as corto hasta un servidor?}

\begin{tmconcepto}{puntos y l\'ineas}
Un grafo es un \textbf{mapa de puntos} (v\'ertices) unidos por l\'ineas (aristas). Sirve para modelar redes, rutas y dependencias.
\end{tmconcepto}

\tmcomo{el algoritmo de Dijkstra}{{Distancia 0 al origen y $\infty$ al resto},{Visita el nodo no visitado m\'as cercano},{Actualiza las distancias de sus vecinos},{Repite hasta visitar todos}}

\subsection{Definici\'on y tipos}
\begin{tmdef}{Grafo}
$G=(V, E)$: un conjunto $V$ de v\'ertices y un conjunto $E$ de aristas. Los v\'ertices son computadoras, usuarios o procesos; las aristas, conexiones, dependencias o relaciones.
\end{tmdef}

\begin{center}
\begin{tikzpicture}[every node/.style={circle,draw=tmAzul,fill=tmAzulClaro,minimum size=0.7cm,font=\sffamily\small},scale=0.9]
  % no dirigido
  \node (n6) at (0,2) {6}; \node (n4) at (1.3,1.2) {4}; \node (n5) at (2.9,1.5) {5};
  \node (n1) at (4.1,0.9) {1}; \node (n3) at (1.0,-0.4) {3}; \node (n2) at (2.8,-0.2) {2};
  \draw[tmAzul] (n6)--(n4)--(n5)--(n1)--(n2)--(n3)--(n4); \draw[tmAzul] (n5)--(n2);
  \node[draw=none,fill=none,font=\sffamily\bfseries\small] at (2,-1.4) {no dirigido};
  % dirigido
  \begin{scope}[xshift=6cm]
  \node (A) at (0,0.8) {A}; \node (B) at (1.3,1.8) {B}; \node (C) at (1.3,-0.2) {C};
  \node (D) at (2.8,-0.2) {D}; \node (E) at (2.8,1.8) {E};
  \foreach \p/\q in {A/B,A/C,B/C,B/E,C/D,D/E}{\draw[-{Stealth},tmAzul] (\p) -- (\q);}
  \node[draw=none,fill=none,font=\sffamily\bfseries\small] at (1.4,-1.2) {dirigido};
  \end{scope}
  % ponderado
  \begin{scope}[xshift=11cm]
  \node (a) at (0,0.8) {A}; \node (b) at (1.2,1.8) {B}; \node (c) at (2.6,1.8) {C};
  \node (d) at (2.6,-0.2) {D}; \node (e) at (1.2,-0.2) {E};
  \draw[tmAzul] (a)--node[draw=none,fill=none,above left,font=\sffamily\scriptsize]{7}(b)
     (b)--node[draw=none,fill=none,above,font=\sffamily\scriptsize]{3}(c)
     (a)--node[draw=none,fill=none,below left,font=\sffamily\scriptsize]{1}(e)
     (b)--node[draw=none,fill=none,left,font=\sffamily\scriptsize]{8}(e)
     (c)--node[draw=none,fill=none,right,font=\sffamily\scriptsize]{6}(d)
     (e)--node[draw=none,fill=none,below,font=\sffamily\scriptsize]{7}(d)
     (c)--node[draw=none,fill=none,above right,font=\sffamily\scriptsize]{2}(e);
  \node[draw=none,fill=none,font=\sffamily\bfseries\small] at (1.4,-1.2) {ponderado};
  \end{scope}
\end{tikzpicture}
\end{center}

\subsection{Grado, caminos y matriz de adyacencia}
\begin{tmdef}{Grado}
En un grafo no dirigido, $\deg(v)$ es el n\'umero de aristas incidentes en $v$. Siempre se cumple $\sum_v\deg(v)=2|E|$.
\end{tmdef}

\begin{tmejemplo}{Grados en una red}
Dispositivos $A, B, C, D$ con $E=\{(A, B),(A, C),(B, C),(C, D)\}$.
\begin{tmpasos}
\item $\deg(A)=2$, $\deg(B)=2$, $\deg(C)=3$, $\deg(D)=1$.
\item Suma: $2+2+3+1=8=2\cdot4$. Correcto.
\item $C$ tiene el mayor grado: se conecta con tres dispositivos.
\end{tmpasos}
\end{tmejemplo}

\subsection{Camino m\'as corto: Dijkstra}
Con pesos \textbf{no negativos}, Dijkstra calcula caminos de costo m\'inimo. Mejor distancia desde $A$ en el grafo ponderado de arriba:

\begin{center}
\begin{tabular}{|c|c|c|c|c|c|}
\hline
\rowcolor{tmAzulClaro}\textbf{Paso} & \textbf{Visita} & \textbf{$B$} & \textbf{$C$} & \textbf{$D$} & \textbf{$E$}\\ \hline
0 & $A$ & 7 & $\infty$ & $\infty$ & 1\\ \hline
1 & $E$ & 7 & 3 & 8 & \textbf{1}\\ \hline
2 & $C$ & 6 & \textbf{3} & 8 & 1\\ \hline
3 & $B$ & \textbf{6} & 3 & 8 & 1\\ \hline
4 & $D$ & 6 & 3 & \textbf{8} & 1\\ \hline
\end{tabular}
\end{center}
Costo $A\to D$: $8$ (camino $A\to E\to D$).

\begin{tmsistemas}
\textbf{Redes:} topolog\'ia, conectividad, fallos. \textbf{IA:} grafos de conocimiento y recomendaci\'on (usuarios y productos como v\'ertices). \textbf{Software:} un grafo dirigido ac\'iclico (DAG) modela dependencias: instalar $\to$ compilar $\to$ probar, y detecta dependencias circulares. La ruta real que usa una red tambi\'en depende de protocolos, congesti\'on y pol\'iticas.
\end{tmsistemas}

\begin{tmcuidado}Dijkstra \textbf{no funciona con pesos negativos}. Y el camino de costo m\'inimo matem\'atico no siempre es la ruta que usa una red real.\end{tmcuidado}

\begin{tmejercicios}[Ejercicios 2.6 \hfill {\small\mdseries respuestas en p\'ag. \pageref{resp:2-6}}]
\label{ej:2-6}
\begin{tmlista}
\item \tmbasico\ Calcula el grado de cada v\'ertice de $V=\{A, B, C, D\}$ con $E=\{(A, B),(A, C),(B, D),(C, D)\}$.
\item \tmbasico\ Si un grafo no dirigido tiene 7 aristas, \textquestiondown cu\'anto suman sus grados?
\item \tmbasico\ Representa las dependencias: instalar dependencias, compilar, ejecutar pruebas.
\item \tmmedio\ Escribe la matriz de adyacencia del grafo del ejercicio 1.
\item \tmmedio\ Con $E=\{(A, B),(C, D)\}$, \textquestiondown existe camino de $A$ a $D$?
\item \tmmedio\ En el grafo ponderado de la secci\'on, halla el camino m\'as corto de $A$ a $B$.
\item \tmdificil\ \textquestiondown Cu\'antas aristas tiene como m\'aximo un grafo simple no dirigido de 8 v\'ertices?
\item \tmdificil\ Explica la diferencia entre usar un grafo para una red f\'isica y para relaciones entre usuarios.
\end{tmlista}
\end{tmejercicios}

\tmreto{2}{Misi\'on 3 (Enlace Cr\'itico): simula el aislamiento del nodo cr\'itico $D$ desactivando el enlace $C$--$D$.}

\begin{tmresumen}
\begin{itemize}
\item $G=(V, E)$; puede ser dirigido, no dirigido o ponderado.
\item $\sum\deg(v)=2|E|$.
\item Dijkstra da caminos de costo m\'inimo con pesos no negativos.
\end{itemize}
\end{tmresumen}

% =====================================================================
\section{Aplicaciones integradoras}
\label{sec:2-7}
\tmgancho{\textquestiondown C\'omo se combinan conteo, recurrencias, teor\'ia de n\'umeros y grafos para resolver un problema completo?}

\begin{tmconcepto}{un problema, varias herramientas}
Un problema real casi nunca es de un solo tema: se \textbf{modela} con grafos, se \textbf{cuenta} con combinatoria y se \textbf{protege} con aritm\'etica modular.
\end{tmconcepto}

\tmcomo{analizar una red}{{Dibuja el grafo},{Calcula los grados},{Busca caminos entre los nodos},{Prueba qu\'e pasa si falla una arista}}

\begin{center}
\begin{tikzpicture}[every node/.style={draw=tmAzul,fill=tmAzulClaro,rounded corners=3pt,font=\sffamily\scriptsize,inner sep=3pt},scale=0.95]
  \node (fw) at (7,3) {Firewall}; \node (in) at (9.2,3) {Internet};
  \node (ro) at (5,2) {Router}; \node (s1) at (2.5,2) {Switch 1};
  \node (web) at (0.2,3) {WebServer}; \node (ftp) at (0,2) {FTP}; \node (em) at (0.2,1) {Email};
  \node (s2) at (3.3,0.6) {Switch 2};
  \node (h1) at (2,-0.5) {Host 1}; \node (h2) at (3.3,-0.5) {Host 2}; \node (h3) at (4.6,-0.5) {Host 3};
  \node (s3) at (7,1.4) {Switch 3}; \node (h4) at (8.4,0.3) {Host 4}; \node (h5) at (9.6,0.3) {Host 5};
  \node (s4) at (6.2,0) {Switch 4}; \node (h6) at (5.2,-1.1) {Host 6}; \node (h7) at (6.3,-1.1) {Host 7}; \node (h8) at (7.4,-1.1) {Host 8};
  \draw[tmAzul] (fw)--(in) (fw)--(ro) (ro)--(s1) (s1)--(web) (s1)--(ftp) (s1)--(em) (ro)--(s2) (s2)--(h1) (s2)--(h2) (s2)--(h3) (ro)--(s3) (s3)--(h4) (s3)--(h5) (s3)--(s4) (s4)--(h6) (s4)--(h7) (s4)--(h8);
\end{tikzpicture}
\end{center}

\begin{center}
\begin{tabular}{|p{3.2cm}|p{3.6cm}|p{6cm}|}
\hline
\rowcolor{tmAzulClaro}\textbf{Caso} & \textbf{Conceptos} & \textbf{Idea}\\ \hline
Redes inform\'aticas & grafos, caminos, conteo & modelar dispositivos y conexiones para analizar rutas y alternativas\\ \hline
Inteligencia artificial & sucesiones, recurrencias, conteo & \'arboles de b\'usqueda: estimar estados y costos\\ \hline
Blockchain & teor\'ia de n\'umeros, aritm\'etica modular & clave p\'ublica; el hash vincula bloques y detecta alteraciones\\ \hline
Ciberseguridad & conteo, divisibilidad, mod & espacios de b\'usqueda y criptograf\'ia; importan implementaci\'on y gesti\'on de claves\\ \hline
Hardware digital & aritm\'etica modular, sucesiones & contadores c\'iclicos e \'indices circulares\\ \hline
\end{tabular}
\end{center}

\subsection{Caso integrador: an\'alisis de una red de servidores}
Cuatro servidores $A, B, C, D$ con $E=\{(A, B),(A, C),(B, C),(C, D)\}$, conexiones bidireccionales del mismo costo.

\begin{tmejemplo}{Cuatro preguntas}
\begin{tmpasos}
\item \textbf{\textquestiondown Cu\'antas conexiones?} $|E|=4$.
\item \textbf{\textquestiondown Qui\'en tiene m\'as conexiones directas?} $C$, con grado 3.
\item \textbf{\textquestiondown Hay camino de $A$ a $D$?} S\'i: $A\to C\to D$ (tambi\'en $A\to B\to C\to D$).
\item \textbf{\textquestiondown Qu\'e pasa si falla $C$--$D$?} $D$ queda aislado: es una conexi\'on cr\'itica.
\end{tmpasos}
\end{tmejemplo}

\begin{tmsistemas}
Un modelo matem\'atico sencillo detecta puntos cr\'iticos y fundamenta decisiones. En una red real hay que considerar tambi\'en la redundancia f\'isica, los protocolos, la capacidad y las condiciones de operaci\'on.
\end{tmsistemas}

\begin{tmejercicios}[Ejercicios 2.7 \hfill {\small\mdseries respuestas en p\'ag. \pageref{resp:2-7}}]
\label{ej:2-7}
\begin{tmlista}
\item \tmbasico\ En el caso integrador, \textquestiondown cu\'antas aristas habr\'ia que agregar para que $D$ tenga dos conexiones?
\item \tmmedio\ Una red tiene 4 servidores y se desea una conexi\'on entre cada par. \textquestiondown Cu\'antos cables hacen falta?
\item \tmmedio\ Si cada cable puede fallar o no (2 estados), \textquestiondown cu\'antos estados distintos tiene una red con 6 cables?
\item \tmdificil\ Propone una situaci\'on donde combines un grafo con aritm\'etica modular (por ejemplo, rutas circulares).
\end{tmlista}
\end{tmejercicios}

\tmreto{2}{Resuelve el caso integrador en el laboratorio: identifica el enlace cr\'itico de la red.}

\begin{tmresumen}
\begin{itemize}
\item Un problema real mezcla varias herramientas: modelo con grafos, conteo para los casos, modular para ciclos.
\item Detectar conexiones cr\'iticas es una aplicaci\'on directa de los grafos.
\item El modelo matem\'atico simplifica; la realidad agrega protocolos, capacidad y redundancia.
\end{itemize}
\end{tmresumen}

% =====================================================================
\clearpage
\section{Actividades de cierre del cap\'itulo}
\label{sec:2-8}

\subsection{Mapa mental}
\tmmapa{Matem\'atica discreta}{{Conteo\\$N=n_1n_2\cdots n_k$},{Combinatoria\\$P(n, r)$ y $\binom nr$},{Recurrencias\\$2T(n/2)+cn$},{Primos\\Euclides y RSA},{Modular\\$a\equiv b\pmod m$},{Grafos\\grado y Dijkstra}}

\subsection{Hoja de f\'ormulas}
\begin{center}
\begin{tabular}{|p{4.4cm}|p{8.2cm}|}
\hline
\rowcolor{tmAzulClaro}\textbf{Tema} & \textbf{F\'ormula}\\ \hline
Aditivo / multiplicativo & $n_1+\cdots+n_k$ \quad / \quad $n_1\cdots n_k$\\ \hline
Permutaci\'on / combinaci\'on & $\frac{n!}{(n-r)!}$ \quad / \quad $\frac{n!}{r!(n-r)!}$\\ \hline
Aritm\'etica / geom\'etrica & $a_n=a_1+(n-1)d$ \quad / \quad $a_n=a_1r^{n-1}$\\ \hline
Divide y vencer\'as & $T(n)=2T(n/2)+cn\Rightarrow\Theta(n\log n)$\\ \hline
MCD y MCM & $\operatorname{MCD}(a, b)\operatorname{MCM}(a, b)=ab$\\ \hline
RSA de juguete & $c=m^e\bmod n$, $m=c^d\bmod n$\\ \hline
Grafos & $\sum\deg(v)=2|E|$\\ \hline
\end{tabular}
\end{center}

\subsection{Actividades}
\begin{tmlista}
\item \textbf{Conteo.} Calcula cu\'antas configuraciones existen con 4 navegadores, 2 sistemas operativos y 3 perfiles, todas v\'alidas.
\item \textbf{Combinatoria.} \textquestiondown Cu\'antos grupos de 3 m\'odulos salen de un sistema de 7?
\item \textbf{Recurrencias.} Analiza $T(n)=2T(n/2)+n$ y di qu\'e algoritmo puede tener ese patr\'on.
\item \textbf{Teor\'ia de n\'umeros.} Calcula $\operatorname{MCD}(252, 105)$ con Euclides.
\item \textbf{Aritm\'etica modular.} Un contador de m\'odulo 16 est\'a en 13. \textquestiondown Qu\'e estado alcanza tras 7 incrementos?
\item \textbf{Grafos.} Representa una red de 6 dispositivos, halla sus grados y determina si todos est\'an conectados.
\item \textbf{Ciberseguridad.} Explica c\'omo el tama\~no de un espacio de b\'usqueda ayuda a evaluar una pol\'itica de contrase\~nas y por qu\'e no basta.
\item \textbf{Integraci\'on.} Propone una situaci\'on de Ingenier\'ia de Sistemas con al menos dos conceptos del cap\'itulo, formula el problema, resu\'elvelo y justifica las herramientas.
\end{tmlista}

\begin{tmsistemas}
\textbf{Bibliograf\'ia inicial:} ``Matem'aticas para Ciencias de Computaci'on'' (Lehman, Leighton y Meyer, MIT OpenCourseWare); ``Matem'aticas Discretas: una introducci'on abierta'' (Levin); recursos criptogr\'aficos del NIST.
\end{tmsistemas}

\tmreto{2}{Cierre del cap\'itulo: completa las tres misiones del laboratorio (espacio de pruebas, contador modular y enlace cr\'itico) y acredita tus TortuPuntos.}

\begin{tmquiz}[Autoevaluaci\'on 2]
\label{quiz:2}
\tmpreg{Con 3 navegadores, 2 sistemas y 4 perfiles hay}{9}{24}{12}{6}
\tmpreg{$P(5, 3)$ vale}{10}{60}{120}{15}
\tmpreg{$\binom83$ vale}{56}{24}{336}{8}
\tmpreg{Con $a_n=a_{n-1}+3$ y $a_1=2$, $a_4$ vale}{8}{11}{14}{9}
\tmpreg{$T(n)=2T(n/2)+cn$ es}{$\Theta(n)$}{$\Theta(n\log n)$}{$\Theta(n^2)$}{$\Theta(2^n)$}
\tmpreg{$\operatorname{MCD}(48, 18)$ vale}{3}{6}{12}{9}
\tmpreg{$17\bmod5$ vale}{2}{3}{1}{0}
\tmpreg{$(7+1)\bmod8$ vale}{8}{1}{0}{7}
\tmpreg{En $E=\{AB, AC, BC, CD\}$, el grado de $C$ es}{1}{2}{3}{4}
\tmpreg{Dijkstra exige que los pesos sean}{negativos}{no negativos}{todos iguales}{enteros}
\end{tmquiz}
